\subsection{Theta-分层}%
\label{sub:Theta-stratifications}

这里使用的 \(\Theta\)-分层与数值不变量的框架发展于~\cite{HLinstability}。

设 \(\mf{X}\) 为代数叠，并考虑 \(\Theta\coloneqq[\A^1/\G_m]\)，其中它在点上的作用为 \(\lambda\cdot z=\lambda z\)；等价地，坐标函数 \(t\) 的权重为 \(-1\)。假设下面出现的映射叠都是代数的。定义滤过对象的叠 \(\ms{Filt}(\mf{X})\coloneqq\ms{Map}(\Theta,\mf{X})\) 与分次对象的叠 \(\ms{Grad}(\mf{X})\coloneqq\ms{Map}(B\G_m,\mf{X})\)。于是有典范态射
\[
    \ev_1\colon\ms{Filt}(\mf{X})\to\mf{X},
    \qquad
    \on{gr}\colon\ms{Filt}(\mf{X})\to\ms{Grad}(\mf{X}),
    \qquad
    \sigma\colon\ms{Grad}(\mf{X})\to\ms{Filt}(\mf{X}).
\]
这里 \(\ev_1\) 是到 \(1\in\A^1\) 的限制，\(\on{gr}\) 由包含 \(B\G_m=\{0\}/\G_m\hookrightarrow\Theta\) 诱导，\(\sigma\) 由投影 \(\Theta\to B\G_m\) 诱导。特别地，\(\on{gr}\circ\sigma=\on{id}\)。我们还记 \(\on{tot}=\ev_1\circ\sigma\colon\ms{Grad}(\mf{X})\to\mf{X}\)。

\begin{defn}
    \textit{\(\Theta\)-层（stratum）} 是代数叠 \(\mf{Y}\) 中连通分支之并 \(S\subset\ms{Filt}(\mf{Y})\)，使得限制
    \[
        \ev_1|_S\colon S\longrightarrow\mf{Y}
    \]
    是闭浸入。

    \textit{\(\Theta\)-分层} 由 \(\mf{X}\) 上的以下数据组成
    \begin{enumerate}
        \item 一个全序集 \(I\)，带有一个最小元 \(0\)；
        \item 一族递增的开子叠 \(\mf{X}_{\leq\alpha}\subset\mf{X}\)，以 \(\alpha\in I\) 为指标，满足
        \[
            \mf{X}_{\leq\alpha}\subset\mf{X}_{\leq\beta}
            \quad\text{for }\alpha\leq\beta,
            \qquad
            \bigcup_{\alpha\in I}\mf{X}_{\leq\alpha}=\mf{X},
        \]
        并且对每个 \(x\in|\mf{X}|\)，集合 \(\{\alpha\in I\mid x\in|\mf{X}_{\leq\alpha}|\}\) 都有最小元；
        \item 对每个 \(\alpha>0\)，一个 \(\Theta\)-层 \(S_\alpha\subset\ms{Filt}(\mf{X}_{\leq\alpha})\)，使得
        \[
            \mf{X}_{<\alpha}
            \coloneqq
            \bigcup_{\beta<\alpha}\mf{X}_{\leq\beta}
            =
            \mf{X}_{\leq\alpha}\setminus\ev_1(S_\alpha).
        \]
    \end{enumerate}
    半稳定轨迹为 \(\mf{X}^{\mr{ss}}\coloneqq\mf{X}_{\leq0}\)，该分层给出一个无交分解
    \[
        \mf{X}
        =
        \mf{X}^{\mr{ss}}
        \sqcup
        \bigsqcup_{\alpha>0}\ev_1(S_\alpha).
    \]
\end{defn}

考虑叠 \(\ms{Bun}_G(C) = \ms{Map}(C, BG)\)。我们将考虑叠
\[ \ms{Map}(\Theta, \ms{Bun}_G(C)) = \ms{Map}(\Theta \times C, BG). \]
任何这样的映射 \(f\) 都等价于一个 \(\G_m\)-等变主 \(G\)-丛 \(\mc{E}\)，定义在 \(\A^1\times C\) 上。当 \(G=\on{GL}_n\) 时，这等价于一个 \(\G_m\)-等变向量丛（定义在 \(\A^1\times C\) 上），或者等价地，一个有限生成的分次局部自由 \(\mc{O}_C[t]\)-模
\[
    \mc{E}=\bigoplus_{w\in\Z}\mc{E}_w.
\]
按上述约定，坐标函数 \(t\) 的权重为 \(-1\)，故乘以 \(t\) 给出映射
\[
    t\colon\mc{E}_w\longrightarrow\mc{E}_{w-1}.
\]
由于 \(\mc{E}\) 在 \(\mc{O}_C[t]\) 上局部自由，乘以 \(t\) 是单射。

\begin{rmk}
    考虑向量丛 \(E=\mc{E}|_{\{1\}\times C}\)。由于 \(\G_m\) 在 \(\G_m\times C\) 上自由作用，等变下降给出
    \[
        \mc{E}|_{\G_m\times C}\simeq P_C^*E,
    \]
    其中 \(P_C\colon\G_m\times C\to C\) 是投影。等价地，
    \[
        E\simeq\big(\mc{E}[t^{-1}]\big)_0.
    \]
    定义
    \[
        F^wE\coloneqq t^w\mc{E}_w\subset\big(\mc{E}[t^{-1}]\big)_0=E.
    \]
    由 \(t\mc{E}_{w+1}\subset\mc{E}_w\) 可知 \(F^{w+1}E\subset F^wE\)。\(\mc{E}\) 的有限生成性与局部自由性蕴含这是一个有限、穷尽且分离的子丛降滤过，它可以写成
    \[
        0=F^{b+1}E\subset F^bE\subset\cdots\subset F^aE=E.
    \]

    反之，设 \(E\) 为带有限降滤过 \(F^\bullet E\)（由子丛给出）的向量丛，且商 \(F^wE/F^{w+1}E\) 局部自由。Rees 构造是如下分次 \(\mc{O}_C[t]\)-模
    \[
        \ms{Rees}(E,F^\bullet)
        \coloneqq
        \bigoplus_{w\in\Z}t^{-w}F^wE
        \subset
        E\otimes_{\C}\C[t,t^{-1}].
    \]
    它在 \(\mc{O}_C[t]\) 上局部自由，且其在 \(t=1\) 与 \(t=0\) 处的纤维满足
    \[
        \ms{Rees}(E,F^\bullet)|_{t=1}\simeq E,
        \qquad
        \ms{Rees}(E,F^\bullet)|_{t=0}
        \simeq
        \bigoplus_w F^wE/F^{w+1}E.
    \]
\end{rmk}

因此，从 \(\Theta\) 到向量丛叠的映射等价于向量丛的一个有限 \(\Z\)-加权滤过，该滤过由子丛给出且相伴分次局部自由。


\subsection{非阿贝尔局部化}%
\label{sub:Non-abelian localization}

我们从经典情形开始。设 \(X\) 为带有 \(\G_m\) 作用的光滑射影簇，设 \(F\) 为 \(\G_m\)-等变完美复形，并记
\[
    X^{\G_m}=\bigsqcup_\alpha Z_\alpha
\]
为不动点轨迹到连通分支的分解。在将表示环 \(R(\G_m)\) 局部化（使得下面的等变欧拉类可逆）之后，我们有
\[
    \chi_{\G_m}(X,F)
    =
    \sum_\alpha
    \chi_{\G_m}\ab(
        Z_\alpha,
        F|_{Z_\alpha}\otimes e(N_{Z_\alpha/X})^{-1}
    ),
\]
其中
\[
    e(N_{Z_\alpha/X})
    =
    \lambda_{-1}(N_{Z_\alpha/X}^\vee)
    =
    \sum_{i\geq0}(-1)^i[\Lambda^iN_{Z_\alpha/X}^\vee].
\]

现在考虑叠版本。设
\[
    \mf{X}
    =
    \mf{X}^{\mr{ss}}
    \sqcup
    \bigsqcup_\alpha\ev_1(S_\alpha)
\]
为一个 \(\Theta\)-分层。由于 \(\ev_1|_{S_\alpha}\) 是闭浸入，可将 \(S_\alpha\) 与其在 \(\mf{X}\) 中的像等同。该层的中心为
\[
    Z_\alpha\coloneqq\sigma^{-1}(S_\alpha)\subset\ms{Grad}(\mf{X}),
\]
其中 \(\sigma\colon\ms{Grad}(\mf{X})\to\ms{Filt}(\mf{X})\) 由投影 \(\Theta\to B\G_m\) 诱导。

\begin{thm}[虚非阿贝尔局部化 {\cite[Theorem~5.1]{HLlocalization}}]
    设 \(\mf{X}\) 为拟光滑导出叠，带有 \(\Theta\)-分层
    \[
        \mf{X}
        =
        \mf{X}^{\mr{ss}}
        \sqcup
        \bigsqcup_\alpha S_\alpha,
    \]
    其中我们将每个层与其在 \(\ev_1\) 下的像等同，并设 \(Z_\alpha\) 为其中心。将余切复形的限制按 \(\G_m\)-权重分解为
    \[
        \L_{\mf{X}}|_{Z_\alpha}
        \simeq
        \L_\alpha^-
        \oplus
        \L_\alpha^0
        \oplus
        \L_\alpha^+,
    \]
    其中 \(\L_\alpha^-\)、\(\L_\alpha^0\) 与 \(\L_\alpha^+\) 的同调分别集中在负权重、零权重与正权重。定义
    \[
        \E_\alpha
        \coloneqq
        \on{Sym}\ab(
            \L_\alpha^-
            \oplus
            (\L_\alpha^+)^\vee
        )
        \otimes
        \det(\L_\alpha^+)^\vee
        [-\on{rk}(\L_\alpha^+)]
        \in
        \ms{QCoh}(Z_\alpha).
    \]
    这里 \(\det(\L_\alpha^+)\) 与 \(\on{rk}(\L_\alpha^+)\) 分别表示完美复形 \(\L_\alpha^+\) 的行列式线丛与虚秩。我们称一个向量空间的复形是有限维的，如果其所有上同调群的直和是有限维的。
    设 \(F\in\ms{Perf}(\mf{X})\)。假设
    \[
        R\Gamma\ab(
            Z_\alpha,
            F|_{Z_\alpha}\otimes\E_\alpha
        )
    \]
    对每个 \(\alpha\) 都是有限维的，且对除有限多个 \(\alpha\) 外都是零调的。则 \(R\Gamma(\mf{X},F)\) 有限维当且仅当 \(R\Gamma(\mf{X}^{\mr{ss}},F)\) 有限维，并且
    \[
        \chi(\mf{X},F)
        =
        \chi(\mf{X}^{\mr{ss}},F)
        +
        \sum_\alpha
        \chi\ab(
            Z_\alpha,
            F|_{Z_\alpha}\otimes\E_\alpha
        ).
    \]
    这里拟光滑是指 \(\L_{\mf{X}}\) 是完美复形且其 Tor-振幅含于 \([-1,1]\)。
\end{thm}


\subsection{几何不变量理论}%
\label{sub:Geometric invariant theory}

关于几何不变量理论，参见~\cite{MFK1994}。下面的向量丛例子用到 Harder--Narasimhan 滤过~\cite{HN1975}。

设 \(G\) 为约化群，\(X\) 为配备 \(G\)-等变线丛 \(\mc{L}\) 的射影概形。半稳定轨迹可以用 Hilbert-Mumford 判别法定义：它由所有满足如下条件的点组成：对所有余特征标
\(\lambda \colon \G_m \to G\)，点 
\[ x_0 = \lim_{t \to 0} \lambda(t) x \]
有定义，并可定义量
\[ \mu^L(x, \lambda) \coloneqq \frac{-\on{wt}_{x_0}(\mc{L})}{\norm{\lambda}}. \]
若 \(\mu^L(x, \lambda) \leq 0\) 对 \(G\) 的所有余特征标成立，则称该点是半稳定的。稳定轨迹由所有满足不等号严格且稳定化子有限的点组成。

\begin{thm}[{\cite{MFK1994}}]
    稳定轨迹与半稳定轨迹在 \(X\) 中是开的。此外，\(X^{\on{ss}}/G\) 是好商（good quotient），而 \(X^s/G\) 是几何商。  
\end{thm}

\begin{exm}
    考虑模叠 \(\ms{Bun}_{r,d}\)，即秩 \(r\)、次数 \(d\) 的向量丛的模叠，其中曲线的亏格至少为 \(2\)。回顾半稳定轨迹被定义为由斜率半稳定（slope semistable）的向量丛组成。事实上，这可以通过 GIT 得到。

    固定一个丰富线丛 \(\mc{O}(1)\)，定义在 \(C\) 上。则存在某个 \(n_0\)，使得对所有 \(n \geq n_0\) 与 \(E \in \ms{Bun}_{r,d}\)，\(E(n)\) 都是整体生成的。因此可以考虑
    \[ H^0(C, E(n)) \twoheadrightarrow E(n). \]
    若令 \(V = \C^N\)，其中 \(N = h^0(E(n))\)，则有满射
    \[ V \otimes \mc{O}_C(-n) \twoheadrightarrow E. \]
    因此可将 \(E\) 视为某个 Quot 概形中的一个成员。在 Quot 概形上，可以考虑泛族
    \begin{equation*}
    \begin{tikzcd}
        V \otimes \mc{O}_C(-n) \ar[twoheadrightarrow]{r} \ar{d} & \mc{E} \\
        \ms{Quot}_{V \otimes \mc{O}_C(-n), p_E} \times C \ar{r}{p} & \ms{Quot}_{V \otimes \mc{O}_C(-n)}.
    \end{tikzcd}
    \end{equation*}
    现在定义线丛
    \[ \mc{L} \coloneqq \det R^0 p_* (\mc{E}(m)) \]
    其中 \(m \gg 0\)。我们也可以考虑 \(\mr{SL}(V)\) 在 Quot 概形上的作用。对某个余特征标 \(\lambda\)，可分解为
    \[ V = V_1 \oplus V_2 \]
    二者的权重分别为 \(s\) 与 \(N-s\)。
    然后考虑
    \[ V_1 \otimes \mc{O}_C(-n) \subset V \otimes \mc{O}_C(-n) \twoheadrightarrow E \]
    以及滤过
    \[ 0 \to F^1 E \to E \to E / F^1 E \to 0 \]
    它由上述包含诱导。于是可得
    \[ \lim_{t \to 0} \lambda(t) E = F^1 E \oplus E/F^1 E. \]
    进而计算限制 \(\mc{L}_{|\on{gr}_F E}\)，得到空间 \(H^0(C, F^1 E(m))\) 与 \(H^0(C, E(m)/F^1 E(m))\)，故 \(\mc{L}\) 的权重为
    \[ (N-s) p_{F^1 E}(m) - s p_{E/F^1 E}(m), \]
    这就恢复了 GIT。
\end{exm}

现在考虑不稳定轨迹。设 \(E\) 为 \(C\) 上的向量丛。则存在唯一的滤过
\[ 0 \subset F^{\ell} E \subseteq \cdots \subseteq F^0 E = E \]
使得
\begin{enumerate}
    \item 所有 \(\on{gr}_F^i E\) 都是半稳定的；
    \item \(\mu(\on{gr}_F^i E) > \mu(\on{gr}_F^{i-1} E)\) 对所有 \(i\) 成立。 
\end{enumerate}
\begin{rmk}
    从 GIT 的观点看，其实应取 \(\mu\) 为约化希尔伯特多项式。
\end{rmk}

\subsection{Theta-分层的构造}%
\label{sub:Construction of Theta-stratifications}

设 \(\mf{X}\) 为一个合理的代数叠。\footnote{这里“合理”是指诸如局部有限表示且对角线为 qcqs 之类。} 现在考虑某个族 \(\mc{E}\)，定义在 \(\A^1 \times C\) 上。若限制到中心纤维，可见 \(\G_m\) 作用在 \(\on{gr}_F E\) 上，即我们有映射 \(\G_m \to \Aut \on{gr}_F E\)，由此可以算出权重。

注意到态射 \(\sigma \colon \ms{Grad}(\mf{X}) \to \ms{Filt}(\mf{X})\) 是一个 \(\Theta\)-收缩（contraction），因此
\[ \pi_0(\ms{Filt}(\mf{X})) = \pi_0(\ms{Grad}(\mf{X})). \]

\begin{exm}
    设 \(\mf{X} = BG\)。则有
    \[\ms{Filt}(\mf{X}) = \bigsqcup_{[\lambda]} B P_{\lambda}, \]
    其中 \(P_{\lambda}\) 是对应于余特征标 \(\lambda\) 的抛物子群，而且我们总可以在单个 Weyl 室中选取代表元。此外还有
    \[ \ms{Grad}(\mf{X}) = \bigsqcup_{\lambda \in \mathbb{X}_*(T)/W} B L_{\lambda}, \]
    其中 \(L_{\lambda}\) 是对应于 \(\lambda\) 的 Levi 子群。
\end{exm}

\begin{defn}
    \textit{\(\Theta\)-层} 是 \(\mf{X}\) 中连通分支之并 \(S \subseteq \ms{Filt}(\mf{X})\)，使得
    \[ \ev_1 \colon S \to \mf{X} \]
    是闭浸入。
\end{defn}

粗略地说，闭浸入是单的，故 \(S\) 中的一个对象由 \(\mf{X}\) 中的一个对象连同一个“典范滤过”组成。

\begin{defn}
    \textit{\(\Theta\)-分层} 是一族递增的开子叠
    \[ \mf{X}_{\leq c} \subseteq \mf{X}, \qquad c \in \Gamma, \]
    其中 \(\Gamma\) 是某个全序集，再配上 \(\Theta\)-层
    \[ S_c \subseteq \ms{Filt}(\mf{X}_{\leq c}), \]
    要求满足
    \[ \mf{X}_{\leq c} = \ev_1(S_c) \sqcup \bigcup_{c' < c} \mf{X}_{\leq c'}. \]
    特别地，这意味着
    \[ \mf{X} = \mf{X}^{\mr{ss}} \sqcup \bigsqcup_{c > 0} \ev_1(S_c). \]
\end{defn}

与 Hilbert-Mumford 判别法类似，我们将用数值不变量来构造 \(\Theta\)-分层。
\begin{defn}
    一个 \textit{数值不变量} \(\mu\)，取值于全序集 \(\Gamma\)，由如下指派组成：指派一个尺度不变函数
    \[ \mu \colon \R^q \setminus \ab\{0\} \to \Gamma \]
    其中 \(p \in \mf{X}(\C)\) 任意，\(\gamma \colon \G_m^q \to \Aut(p)\) 为任意非退化同态。这些指派要求满足
    \begin{enumerate}
        \item 该指派在域扩张下不变。由于我们在 \(\C\) 上工作，这一点不再详述；
        \item 若给出 \(\G_m^q \hookrightarrow \G_m^n \hookrightarrow \Aut(p)\)，则有分解 \(\R^q \setminus \ab\{0\} \to \R^n \setminus \{0\} \to \Gamma\)；
        \item 若给出某个 \(\xi \colon S \to \mf{X}\) 以及核有限的 \(\gamma \colon \G_{m, S}^q \to \Aut(\xi)\)，则函数 \(\mu_{\kappa(s)}\) 在 \(S\) 上局部常值，其中 \(s\) 取遍 \(S\) 的点。
    \end{enumerate}
\end{defn}

我们给出从数值不变量 \(\mu\) 出发的一个朴素构造。定义 \(x \in \mf{X}\) 为半稳定的，如果 \(\mu(f) \leq 0\) 对所有非退化 \(f \in \ms{Filt}(\mf{X})\)（满足 \(f(1) = x\)）成立。若 \(x\) 不稳定，则定义其不稳定水平为
\[ M^{\mu}(x) \coloneqq \sup \ab\{ \mu(f) \mid f \in \ms{Filt}(\mf{X}), \ev_1(f) = x\}. \]
若某个滤过 \(f\) 达到该上确界，则称其为“Harder-Narasimhan 滤过”。

\begin{defn}
    称数值不变量 \(\mu\) 是 \textit{标准的}，如果
    \begin{enumerate}
        \item \(\mu(x)\) 与 \(\mu(-x)\) 不同时为正（\(x \in \R^q \setminus \ab\{0\}\)）；
        \item 每个函数 \(\mu_{\gamma}\) 都是严格拟凹的，即对所有 \(\delta \in [0,1]\) 以及线性无关的 \(x_0, x_1 \in \R^q\)（满足 \(\mu(x_0), \mu(x_1) \geq 0\)），有 
            \[\mu_{\gamma}(\delta x_0 + (1-\delta)x_1) \geq \max\ab\{\mu_{\gamma}(x_0), \mu_{\gamma}(x_1)\}.\] 若其中有一个值严格为正，则要求不等式严格。
    \end{enumerate}
\end{defn}

下一个条件是某种有理性。我们考虑 \(\mathbb{X}_*(T)_{\R}\)，目标是找到有理点。这就引出条件 (R)：若 \(\mu_{\gamma}\) 对某个 \(\gamma\) 取到正值，则它在某个有理点处取到最大值。

\begin{thm}[{\cite{HLinstability}}]
    设 \(\mf{X}\) 为叠，\(\mu\) 为满足条件 (R) 的标准数值不变量。则 \(\mu\) 定义一个 \(\Theta\)-分层当且仅当它满足以下条件：
    \begin{itemize}
        \item 对任意 DVR \(R\)，设其分式域为 \(K\)、剩余域为 \(k\)，则 \(\mu(f_K^{\mr{HN}}) \leq M^{\mu}(f_k)\)（HN-特殊化）；
        \item 对任意从有限型概形出发的映射 \(B \colon T \to \mf{X}\)，存在拟紧开子叠 \(\mf{X}' \subseteq \mf{X}\)，使得对任意有限型点 \(p \in T(k)\) 及滤过 \(f\)（\(B(p)\) 的滤过），若 \(\mu(f) > 0\)，则存在另一个滤过 \(f'\)（同为 \(B(p)\) 的滤过）满足 \(\mu(f') \geq \mu(f)\) 且 \(f'(0) \in \mf{X}'\)（HN-有界性）。
    \end{itemize}
\end{thm}

这里的要点是：标准性与条件 (R) 连同 HN-有界性蕴含存在性，而 HN-特殊化蕴含唯一性。



